\setcounter{section}{2}

\Needspace{5\baselineskip}
\hypertarget{encoderdecoder-architecture}{%
\section{Encoder--Decoder Architecture}\label{encoderdecoder-architecture}}

The encoder--decoder architecture aims to convert an input sequence (for example, an English sentence) (x₁, x₂, \ldots, x\_T) into an output sequence (for example, a Chinese sentence) (y₁, y₂, \ldots, y\_T'). The sequences can differ in length.

Encoder: reads and ``understands'' the entire input sequence, compressing it into a fixed-dimensional context vector c intended to contain its complete semantic information.

Decoder: uses context vector c as its initial ``memory'' and ``instructions'' and generates the output sequence step by step.

\Needspace{5\baselineskip}
\hypertarget{i.-encoder}{%
\subsection{I. Encoder}\label{i.-encoder}}

The encoder is an RNN (a standard RNN, LSTM, or GRU) that processes each input token in order.

\Needspace{5\baselineskip}
At every time step t (from 1 to T), it computes:
h\_t\^{}\{enc\} = f\_\{enc\}(x\_t, h\_\{t-1\}\^{}\{enc\})

h\_t\^{}\{enc\}: the encoder hidden state at time t.

f\_\{enc\}: the encoder RNN's function (such as an LSTM cell's computation). It contains learnable parameters W\_\{xh\}\^{}\{enc\}, W\_\{hh\}\^{}\{enc\}, b\_h\^{}\{enc\}, and others, as well as activation functions.

x\_t: the input word-embedding vector at time t.

h\_\{t-1\}\^{}\{enc\}: the preceding hidden state. The initial state h\_0\^{}\{enc\} is usually an all-zero vector.

\Needspace{5\baselineskip}
Generating the context vector:
After processing the complete input, the encoder's final hidden state h\_T\^{}\{enc\} is used as context vector c.
c = h\_T\^{}\{enc\}

This vector c carries all information read and integrated from the input sequence. It initializes the decoder's state.

In more advanced architectures, c can be a pooled result or an attention-weighted combination, but its most classic form is h\_T\^{}\{enc).

\Needspace{5\baselineskip}
\hypertarget{ii.-decoder}{%
\subsection{II. Decoder}\label{ii.-decoder}}

The decoder is another RNN whose goal is to generate the output sequence. Its initial state is initialized with the encoder's context vector.

\Needspace{5\baselineskip}
Mathematically, its computation at every time t' (from 1 to T') has two parts:

\begin{enumerate}
\def\labelenumi{\alph{enumi})}
\tightlist
\item
\Needspace{5\baselineskip}
  Compute the current hidden state:
  s\_\{t'\} = f\_\{dec\}(y\_\{t'-1\}, s\_\{t'-1\})
\end{enumerate}

s\_\{t'\}: the decoder hidden state at time t'.

f\_\{dec\}: the decoder RNN's function, with parameters W\_\{xh\}\^{}\{dec\}, W\_\{hh\}\^{}\{dec\}, b\_h\^{}\{dec\}, and others.

s\_\{t'-1\}: the preceding decoder hidden state. Crucially, the initial state s\_0 is set to context vector c.
s\_0 = c

y\_\{t'-1\}: the output generated by the decoder at the preceding time step. The input policy differs between training and inference (see ``teacher forcing'' below).

\begin{enumerate}
\def\labelenumi{\alph{enumi})}
\setcounter{enumi}{1}
\tightlist
\item
\Needspace{5\baselineskip}
  Compute the current output:
  Hidden state s\_\{t'\} passes through an output layer (usually fully connected + Softmax) to predict the probability distribution of the current output token.
  o\_\{t'\} = g(s\_\{t'\}) = \text{softmax}(W\_\{out\} s\_\{t'\} + b\_\{out\})
  P(y\_\{t'\} \textbar{} y\_\{\textless t'\}, c) = o\_\{t'\}
\end{enumerate}

o\_\{t'\}: a probability-distribution vector whose dimensionality equals the output-vocabulary size. Each value is the probability that the corresponding token is output at time t'. In general, each word is represented by

W\_\{out\}, b\_\{out\}: the output layer's weight and bias parameters.

\Needspace{5\baselineskip}
The decoder generates a sequence autoregressively:

Initial state s\_0 = c; the initial input is usually a special \texttt{\textless{}start\textgreater{}} token.

Compute s\_1 and o\_1; sample from o\_1 or take argmax to obtain the first word y₁.

Feed y₁ into the next time step, compute s\_2 and o\_2, and obtain y₂.

Repeat until a special \texttt{\textless{}end\textgreater{}} token is generated, indicating completion.

\Needspace{5\baselineskip}
\hypertarget{iii.-summary-and-complete-mathematical-process}{%
\subsection{III. Summary and Complete Mathematical Process}\label{iii.-summary-and-complete-mathematical-process}}

\Needspace{5\baselineskip}
Translate input sequence X = (x₁, x₂, \ldots, x\_T) into output sequence Y = (y₁, y₂, \ldots, y\_\{T'\}):

\begin{enumerate}
\def\labelenumi{\arabic{enumi}.}
\item
\Needspace{5\baselineskip}
  Encoder stage:
  h\_t\^{}\{enc\} = f\_\{enc\}(x\_t, h\_\{t-1\}\^{}\{enc\}), \quad \text{for } t = 1, \ldots, T
  c = h\_T\^{}\{enc\}
\item
\Needspace{5\baselineskip}
  Decoder stage:
  s\_0 = c
  s\_\{t'\} = f\_\{dec\}(y\_\{t'-1\}, s\_\{t'-1\}), \quad \text{for } t' = 1, \ldots, T'
  P(y\_\{t'\} \textbar{} y\_\{\textless t'\}, c) = g(s\_\{t'\}) = \text{softmax}(W\_\{out\} s\_\{t'\} + b\_\{out\})
\item
\Needspace{5\baselineskip}
  Training objective and loss:
\Needspace{5\baselineskip}
  Maximize the output sequence's conditional log-likelihood:
  Maximize log P(Y\textbar X), equivalently sum(log P(y\_\{t'\} \textbar{} y\_\{\textless t'\}). The term -log P(y\_\{t'\} \textbar{} y\_\{\textless t'\}) is cross-entropy loss, pushing the model toward the true one-hot distribution.
\item
  An important technique: teacher forcing
\end{enumerate}

\Needspace{5\baselineskip}
During training, there is an important choice for the input y\_\{t'-1\} at each decoder step:

Without teacher forcing: use the decoder's own preceding prediction as input. Errors accumulate and training may become unstable.

With teacher forcing: regardless of the preceding prediction, forcibly use the previous ground-truth word y\_\{t'-1\}\^{}\{true\} as input.

s\_\{t'\} = f\_\{dec\}(y\_\{t'-1\}\^{}\{true\}, s\_\{t'-1\})

This greatly accelerates convergence by giving the decoder correct context at every step. However, it can create a training--inference mismatch known as exposure bias, because ground truth is unavailable during inference. A curriculum-learning strategy is commonly used to gradually reduce teacher forcing later in training.

At inference time, no ground truth is available, so only the decoder's own preceding prediction can serve as the current input.

\FloatBarrier
\Needspace{5\baselineskip}
\hypertarget{references-19}{%
\subsection{References}\label{references-19}}

\vspace{0.25em}
\begin{itemize}
\tightlist
\item
  Zhang, A., Lipton, Z. C., Li, M., \& Smola, A. J. (2023). \href{https://D2L.ai}{Dive into Deep Learning}. Cambridge University Press.
\item
  Bengio, S., Vinyals, O., Jaitly, N., \& Shazeer, N. (2015). \href{https://proceedings.neurips.cc/paper/2015/hash/e995f98d56967d946471af29d7bf99f1-Abstract.html}{Scheduled Sampling for Sequence Prediction with Recurrent Neural Networks}. NeurIPS 2015.
\end{itemize}

\clearpage
